\documentclass[11pt]{article}
\usepackage[margin=1.02in]{geometry}
\usepackage{amsmath,amssymb,amsthm,mathtools,booktabs,hyperref,microtype}
\hypersetup{colorlinks=true,linkcolor=blue,urlcolor=blue}
\newtheorem{theorem}{Theorem}
\newtheorem{lemma}{Lemma}
\newtheorem{corollary}{Corollary}
\newcommand{\E}{\mathbb E}
\newcommand{\Tr}{\operatorname{Tr}}
\newcommand{\id}{\mathbf 1}
\newcommand{\HH}{\mathsf H}
\title{\bf Sharpness of the Operational SLD Tensor Transfer:\\ Ternary Exponential Hiding and Binary No-Loss}
\author{Devin J. Wright\\\small NIGHTFALL research program; AI-assisted construction and audit}
\date{11 November 2026 \quad\small Research addendum / no external referee or priority certification}
\begin{document}
\maketitle
\begin{abstract}
We resolve the quantitative sharpness question for the $2^{1-d}$ factor in the operational symmetric logarithmic derivative (SLD) transfer of \emph{Universal Observability}. The coefficient is exactly optimal even after restricting to legitimate independent classical inputs and locally prepared pure qubit states with strictly positive local contrast Gram operators. A uniform ternary $d$-coordinate source, a bounded real pure degree-$d$ target with variance $1/2$, and a faithful local qubit average yield a closed-form terminal Bayes gain. At a fixed local visibility $a=3/8$, its ratio to $a^d$ times classical influence decays as $2(2/3)^d$, excluding every polynomial-degree improvement of the generic transfer. A complementary Jensen argument proves a factor-one SLD comparison for simple tensor products of Hermitian local contrasts, implying that all binary-input Hoeffding interactions \emph{avoid} the additional $2^{1-d}$ penalty. A separate commuting binary example shows that the remaining factor $a^d$ cannot generally be removed. These are mathematical derivations with an independently represented finite-matrix checker, not claims of historical originality or independent human review.
\end{abstract}

\section{Operational setting}
For a faithful state $b$, let $J_b(H)=(bH+Hb)/2$ and
\[
 Q_b(H):=\Tr\bigl[H J_b^{-1}(H)\bigr]=\sum_{u,v}\frac{2|H_{uv}|^2}{\lambda_u+\lambda_v}
\]
in an eigenbasis of $b$, where $H=H^*$ and $\lambda_u>0$. For independent finite classical sources $X_i\sim\mu_i$ and locally prepared states $\rho_x=\bigotimes_i\rho_{i,x_i}$, put $b_i=\E\rho_{i,X_i}\succ0$, and define on the centered real source space
\[
 T_i g=\E[g(X_i)\rho_{i,X_i}],\qquad K_i=T_i^*J_{b_i}^{-1}T_i,\qquad a_i=\lambda_{\min}(K_i).
\]
For a real target $f$ with centered Hoeffding component $f_S$, let $A_S=(\bigotimes_{i\in S}T_i)f_S$ and $\gamma_S=Q_{b_S}(A_S)$. Operational quantum ANOVA from the companion manuscript identifies $\gamma_S$ as the nonnegative interaction gain, with terminal gain $\Delta_i=\sum_{S\ni i}\gamma_S$.
The preceding manuscript established
\begin{equation}\label{eq:envelope}
\gamma_S\ge 2^{1-|S|}\langle f_S,(\otimes_{i\in S}K_i)f_S\rangle,
\qquad \Delta_i\ge 2^{1-d}a^d I_i(f)
\end{equation}
when $f$ has degree at most $d$ and $a_i\ge a\in(0,1]$.

\section{Physical sharpness and fixed-visibility separation}
\begin{theorem}[Ternary coherent-phase obstruction]\label{thm:ternary}
Let $d\ge1$, and let $X_1,\dots,X_d$ be independent, uniform on $\{0,1,2\}$. Put $\omega=e^{2\pi i/3}$, choose any $p\in(1/2,1)$, write $q=1-p$, $u=\sqrt{pq}$, and encode each symbol $x$ by the pure qubit
\[
 \rho_x=\begin{pmatrix}p&u\omega^{-x}\\u\omega^x&q\end{pmatrix}
 =\bigl(\sqrt p\,|0\rangle+\sqrt q\,\omega^x|1\rangle\bigr)
 \bigl(\sqrt p\,\langle0|+\sqrt q\,\omega^{-x}\langle1|\bigr).
\]
Define $f_d(x)=\Re\omega^{x_1+\cdots+x_d}$. Then $|f_d|\le1$, its Hoeffding degree is exactly $d$, and
\begin{equation}\label{eq:exact}
 \mathrm{Var}(f_d)=I_i(f_d)=\frac12,\qquad
 K_i=2pq\,\id_{L_0^2(\mu_i)},\quad a_i=a=2pq,
 \qquad\Delta_i=\frac{(pq)^d}{p^d+q^d}.
\end{equation}
Every proper register subset has zero Bayes gain, whereas joint observation of all $d$ registers achieves the stated $\Delta_i$. In particular,
\begin{equation}\label{eq:ratio}
 \boxed{\frac{\Delta_i}{a^d I_i(f_d)}=\frac{2^{1-d}}{p^d+q^d}.}
\end{equation}
For fixed $p=3/4$, $a=3/8$ and
\begin{equation}\label{eq:fixed}
 \boxed{\Delta_i=\frac{3^d}{4^d(3^d+1)},\qquad
 \frac{\Delta_i}{a^dI_i(f_d)}=\frac{2^{d+1}}{3^d+1}\sim2(2/3)^d.}
\end{equation}
Consequently no universal polynomial $P$ and positive constant $c$ give $\Delta_i\ge c a^d I_i(f)/P(d)$ for this class of sources and encodings. Moreover, at each fixed $d$, the coefficient $2^{1-d}$ in the first inequality of \eqref{eq:envelope} is the largest possible uniformly across locally prepared finite-alphabet quantum models with $K_i\succ0$.
\end{theorem}
\begin{proof}
The root-of-unity identities $\E\omega^{X_i}=\E\omega^{2X_i}=0$ imply $\E[f_d\mid X_S]=0$ for every proper $S\subsetneq[d]$. Thus $f_d$ is its sole order-$d$ Hoeffding component. Also $\E f_d=0$ and $\E f_d^2=1/2$, so $I_i(f_d)=1/2$.

The local average is $b=\operatorname{diag}(p,q)\succ0$. An orthonormal basis of the real centered ternary source space consists of $g_c(x)=\sqrt2\cos(2\pi x/3)$ and $g_s(x)=\sqrt2\sin(2\pi x/3)$. The contrast map sends them to $u\sigma_x/\sqrt2$ and $u\sigma_y/\sqrt2$, respectively. Since $J_b^{-1}$ multiplies off-diagonal matrix entries by $2/(p+q)=2$, these images are SLD-orthogonal and have squared SLD energy $2u^2$ each. Hence $K=2pq\,\id$, proving the local claim.

Write $B=b^{\otimes d}$. Expand $\rho_x=b+u(\omega^{-x}|0\rangle\langle1|+\omega^x|1\rangle\langle0|)$. Orthogonality of the ternary characters shows that the target-weighted density matrix has exactly two nonzero entries:
\[
 A=\E\bigl[f_d(X)\rho_{X_1}\otimes\cdots\otimes\rho_{X_d}\bigr]
 =\frac{u^d}{2}\bigl(|0^d\rangle\langle1^d|+|1^d\rangle\langle0^d|\bigr).
\]
These extreme off-diagonal matrix units have zero partial trace over any nonempty factor, hence every proper register subset has zero score and zero Bayes gain. The corresponding $B$-eigenvalues are $p^d$ and $q^d$, so Personick's operational SLD formula yields
\[
 G^*([d])=Q_B(A)=2\cdot\frac{u^{2d}}4\cdot\frac{2}{p^d+q^d}=\frac{u^{2d}}{p^d+q^d}.
\]
Therefore all terminal gains equal $G^*([d])$, and \eqref{eq:ratio}--\eqref{eq:fixed} follow by exact algebra. The fixed-$p$ ratio decreases exponentially, ruling out every polynomial replacement. Finally, for fixed $d$, $p\uparrow1$ makes $p^d+q^d\to1$ while each $K_i=2pq\,\id$ remains positive before the limit. Therefore the universal best constant relative to $\otimes_iK_i$ cannot exceed $2^{1-d}$. Combined with \eqref{eq:envelope}, the coefficient is exact. The limiting visibility tends to zero; this does not undermine the fixed-$a=3/8$ exponential separation.
\end{proof}

\section{No extra quantum penalty for simple Hermitian contrasts}
\begin{lemma}[Factor-one SLD supermultiplicativity on Hermitian simple tensors]\label{lem:product}
For faithful finite-dimensional states $b_i$ and arbitrary Hermitian matrices $H_i$,
\begin{equation}\label{eq:simple}
 \boxed{Q_{\otimes_{i=1}^d b_i}\bigl(\otimes_{i=1}^d H_i\bigr)
 \ge\prod_{i=1}^d Q_{b_i}(H_i).}
\end{equation}
The coefficient $1$ is sharp, even for nonzero Hermitian contrasts.
\end{lemma}
\begin{proof}
Choose eigenbases $b_i=\sum_j\lambda_{ij}|j\rangle\langle j|$. Expand both sides using $Q_b(H)=\sum_{j,k}2|H_{jk}|^2/(\lambda_j+\lambda_k)$. For a fixed tuple of \emph{unordered} eigen-index pairs $\{j_i,k_i\}$, Hermiticity ensures that the absolute-squared matrix-element weight is invariant under independently swapping $j_i$ and $k_i$. If $j_i=k_i$, regard the two orientations as identical. Let $\varepsilon_i$ be independent fair swaps and define
\[
 D(\varepsilon)=\prod_{i=1}^d \lambda_{i,j_i(\varepsilon_i)}
 +\prod_{i=1}^d \lambda_{i,k_i(\varepsilon_i)}.
\]
The relevant global-to-local kernel ratio averaged over orientations is
\[
 \E_{\varepsilon}\left[
 \frac{2^{1-d}\prod_i(\lambda_{ij_i}+\lambda_{ik_i})}{D(\varepsilon)}\right].
\]
Now $\E D(\varepsilon)=2^{1-d}\prod_i(\lambda_{ij_i}+\lambda_{ik_i})$. Jensen's inequality for the convex function $1/t$ implies the displayed mean is at least $1$. Sum over unordered pairs with their common nonnegative weights to get \eqref{eq:simple}. Equality occurs when each $b_i$ is maximally mixed (or for diagonal inputs), so the constant is sharp.
\end{proof}

\begin{corollary}[Binary local alphabets incur no extra SLD tensor loss]\label{cor:binary}
Suppose all nontrivial classical input alphabets are binary, with arbitrary nondegenerate biases and arbitrary (possibly noncommuting) faithful-average local quantum encoders. Then for every classical Hoeffding component $f_S$,
\[
 \boxed{\gamma_S\ge\langle f_S,(\otimes_{i\in S}K_i)f_S\rangle
 \ge\bigl(\prod_{i\in S}a_i\bigr)\|f_S\|_{L^2}^2.}
\]
Consequently, if $\deg(f)\le d$ and $a_i\ge a$, then $\Delta_i\ge a^d I_i(f)$ for every $i$.
\end{corollary}
\begin{proof}
Every centered real $L^2$ space on a binary alphabet is one-dimensional. Thus a pure Hoeffding component is a real multiple of the tensor product of its normalized centered single-coordinate functions. Its encoded contrast is a real multiple of a simple tensor of Hermitian matrices. Apply Lemma~\ref{lem:product}, then the scalar local Gram bound. Summing nonnegative interaction gains over $S\ni i$ proves the terminal statement.
\end{proof}

\begin{corollary}[The factor $a^d$ remains necessary for commuting binary encoders]
For each fixed $0<a<1$ there is a family of independent uniform signs, commuting local qubit encoders with $K_i=a$, and the bounded parity target $f(x)=\prod_i x_i$, such that $\mathrm{Var}(f)=I_i(f)=1$ but $\Delta_i=a^d$ for all $i$.
\end{corollary}
\begin{proof}
For sign input $x\in\{-1,1\}$, take $\rho_x=(I+x\sqrt a\,\sigma_z)/2$, with $b=I/2$. The centered local contrast $T(x)=\sqrt a\,\sigma_z/2$ has SLD energy $a$. The full target-weighted density is $T(x)^{\otimes d}$, and the SLD inverse at $b^{\otimes d}$ is scalar multiplication by $2^d$. Thus its SLD energy equals $a^d$. All proper conditional target scores vanish.
\end{proof}

\section{Interpretation, provenance and next questions}
The full SLD-tensor coefficient $2^{1-d}$ is sharp even for \emph{physical} local preparations; the ternary example supplies a positive-visibility, constant-variance, all-degree obstruction to polynomial replacement. The binary result locates a structural boundary: independent local Hermitian tensor contrasts themselves incur no extra quantum penalty; destructive interference among multiple local contrast directions is essential to the ternary loss. Neither result is an unconditional quantum Aaronson--Ambainis inequality with only polynomial degree dependence: the commuting parity example already makes $a^d$ unavoidable at fixed $a<1$.

The quantum gain here is optimal \emph{single-copy scalar Bayes mean-square gain under unrestricted joint measurements} on the designated registers. The theorem does not concern quantum computational query complexity, many-copy asymptotics, or a restriction to local/separable POVMs. The local source/encoder setup is finite and uses faithful average states; the pure conditional qubits in Theorem~\ref{thm:ternary} cause no singularity because the average is faithful.

The proof uses the operational SLD estimator from Personick and the interaction formalism of the accompanying NIGHTFALL manuscript. Relevant prior art in quantum $\chi^2$ divergence and strong-data-processing tensorization includes Cao--Lu (2019) and Hirche--Rouz\'e--Fran\c ca (2022). This dated topical check has not established historical priority of either the explicit trine-phase separation or the simple-tensor Jensen lemma. Independent specialist review is required before a publishability claim. The accompanying code has two independently represented checks (root-of-unity exact arithmetic and density-matrix computation), plus sampled stress tests; it does not prove the universally quantified Jensen result by enumeration.

\begin{thebibliography}{9}\small
\bibitem{personick} S. Personick, Application of quantum estimation theory to analog communication over quantum channels, \emph{IEEE Trans. Inf. Theory} \textbf{17} (1971), 240--246.
\bibitem{cao} Y. Cao and J. Lu, Tensorization of the strong data processing inequality for quantum chi-square divergences, \emph{Quantum} \textbf{3} (2019), 199. \url{https://arxiv.org/abs/1904.06562}.
\bibitem{hirche} C. Hirche, C. Rouz\'e and D. Stilck Fran\c ca, On contraction coefficients, partial orders and approximation of capacities for quantum channels, \emph{Quantum} \textbf{6} (2022), 862.
\bibitem{wright} D. J. Wright, \emph{Universal Observability: Distribution-Free Classical Prediction and Quantum Information Geometry}, NIGHTFALL pre-referee manuscript (11 November 2026), Section 5.
\end{thebibliography}
\end{document}
